% Part: history
% Chapter: biographies 
% Section: kurt-goedel

\documentclass[../../../include/open-logic-section]{subfiles}

\begin{document}

\olfileid{his}{bio}{god}

\olsection{कुर्ट गोडेल}

\olphoto{goedel-kurt}{कुर्ट गोडेल}

कुर्ट गोडेल (\textsc{गर}-डल) यांचा जन्म २८ एप्रिल १९०६ रोजी ऑस्ट्रो-हंगेरी साम्राज्यातील ब्र्युन येथे झाला; ते ठिकाण आता झेक प्रजासत्ताकातील ब्र्नो आहे. लहानपणी ते जिज्ञासू आणि हुशार असल्याने घरचे त्यांना “Der kleine Herr Warum” (लहानसे 'का-महोदय') असे म्हणत. प्राथमिक शाळेपासूनच त्यांची शैक्षणिक कामगिरी उत्कृष्ट होती; सर्वोच्च गुणांपेक्षा कमी गुण त्यांना फक्त गणितात मिळाले. तब्येत बरी नसल्याने गोडेल अनेकदा शाळेत गैरहजर राहत आणि त्यांना शारीरिक शिक्षणातून सूट मिळाली होती. बालपणी त्यांना संधिवातजन्य तापाचे निदान झाले. वैद्यकीय तपासणीत त्याच्या उलट निष्कर्ष आला असतानाही, त्या आजाराचा आपल्या हृदयावर कायमचा परिणाम झाला आहे, असे त्यांना आयुष्यभर वाटत राहिले.

गोडेल यांनी १९२४ मध्ये व्हिएन्ना विद्यापीठात शिक्षण सुरू केले आणि १९२९ मध्ये डॉक्टरेट पूर्ण केली. सुरुवातीला त्यांचा भौतिकशास्त्र शिकण्याचा विचार होता; पण काही अंशी तत्त्वज्ञ रुडॉल्फ कार्नाप यांच्या प्रभावामुळे त्यांची आवड गणिताकडे, विशेषतः तर्कशास्त्राकडे वळली. हान्स हान यांच्या मार्गदर्शनाखाली लिहिलेल्या प्रबंधात त्यांनी समानतेसह प्रथम-क्रम विधेय तर्कशास्त्राचे संपूर्णता प्रमेय सिद्ध केले \citep{Godel1929}. त्यानंतर अवघ्या वर्षभरात त्यांनी आपले सर्वात प्रसिद्ध निष्कर्ष, पहिले आणि दुसरे अपूर्णता प्रमेय, मिळवले (हे \citealt{Godel1931} मध्ये प्रसिद्ध झाले). व्हिएन्नातील वास्तव्यात गोडेल 'व्हिएन्ना मंडळ' या विज्ञानाभिमुख तत्त्वज्ञांच्या समूहात सक्रिय होते. या समूहातील कार्नाप यांच्या कार्यावर गोडेल यांच्या निष्कर्षांचा विशेष प्रभाव पडला.

१९३८ मध्ये गोडेल यांनी अडेल निंबुर्स्की यांच्याशी विवाह केला. त्यांचे पालक या विवाहावर नाराज होते: अडेल त्यांच्यापेक्षा सहा वर्षांनी मोठ्या होत्या, त्यांचा आधी घटस्फोट झाला होता आणि त्या एका नाइटक्लबमध्ये नर्तकी म्हणून काम करत होत्या. मात्र सामाजिक दबावाचा गोडेल यांच्यावर परिणाम झाला नाही; त्यांच्या निधनापर्यंत त्यांचा विवाह आनंदी राहिला, असे स्रोत सांगतो.

१९३८ मध्ये नाझी जर्मनीने ऑस्ट्रियाचा विलय केल्यानंतर गोडेल आणि अडेल अमेरिकेत स्थलांतरित झाले. न्यू जर्सीतील प्रिन्स्टन येथील इन्स्टिट्यूट फॉर ॲडव्हान्स्ड स्टडी येथे गोडेल यांनी पद स्वीकारले. ते अंतर्मुख आणि काहीसे विलक्षण स्वभावाचे असले, तरी प्रिन्स्टनमधील त्यांचा काळ सहकार्यपूर्ण आणि फलदायी ठरला. त्यांनी संच उपपत्ती, तत्त्वज्ञान आणि भौतिकशास्त्र यांवर निबंध प्रकाशित केले. याच संस्थेतील सहकारी अल्बर्ट आइन्स्टाइन यांच्याशी त्यांची विशेष घनिष्ठ मैत्री झाली.

आयुष्याच्या उत्तरार्धात गोडेल यांचे मानसिक आरोग्य खालावले. १९७७ मध्ये त्यांच्या पत्नीला रुग्णालयात दाखल करावे लागल्याने त्या त्यांच्यासाठी जेवण बनवू शकत नव्हत्या. गोडेल यांना आयुष्यभर मानसिक आरोग्याच्या समस्या होत्या; त्यातून त्यांच्या मनात संशय बळावला. आपल्याला विष घातले जाईल या तीव्र भीतीने त्यांनी अन्न घेणे बंद केले. १४ जानेवारी १९७८ रोजी प्रिन्स्टन येथे उपासमारीने त्यांचे निधन झाले, असे या स्रोताचे वर्णन आहे.

\begin{reading}
गोडेल यांचे सविस्तर चरित्र \citet{Dawson1997} येथे पाहा. त्यांच्या चरित्राविषयीचे आणखी लेख आणि तर्कशास्त्र व तत्त्वज्ञानातील योगदानावरील निबंध \citet{Wang1990}, \citet{Baaz2011}, \citet{Takeuti2003} आणि \citet{Sigmund2007} येथे मिळतील.

गोडेल यांचा डॉक्टरेट प्रबंध मूळ जर्मन भाषेत \citep{Godel1929} येथे उपलब्ध आहे. अपूर्णता प्रमेयांचा मूळ लेख \citep{Godel1931} येथे आहे. त्यांचे प्रकाशित व अप्रकाशित सर्व लेखन आणि निवडक पत्रव्यवहार यांची इंग्रजी आवृत्ती त्यांच्या \emph{Collected Papers} \citet{Godel1986,Godel1990} मध्ये उपलब्ध आहे.

गोडेल यांच्या अपूर्णता प्रमेयांच्या सविस्तर विवेचनासाठी \citet{Smith2013} पाहा. या प्रमेयांवरील अनौपचारिक तात्त्विक चर्चेसाठी मार्क लिन्झेनमायर यांचे पॉडकास्ट \citep{Linsenmayer2014} पाहा.
\end{reading}

\end{document}
